% ============================================================================
% Companion Compendium — Encoded Magic-State Preparation
% A gentle introduction to the physics, engineering, and optimisation
% behind the autoresearch-quantum notebook series
% ============================================================================
\documentclass[11pt,a4paper,twoside]{report}

% ── Typography & layout ─────────────────────────────────────────────────────
\usepackage[T1]{fontenc}
\usepackage[utf8]{inputenc}
\usepackage{lmodern}
\usepackage{microtype}
\usepackage[margin=2.5cm,inner=3cm,outer=2cm]{geometry}
\usepackage{parskip}
\usepackage{setspace}
\onehalfspacing
\usepackage{fancyhdr}
\pagestyle{fancy}
\fancyhf{}
\fancyhead[LE]{\small\itshape\leftmark}
\fancyhead[RO]{\small\itshape\rightmark}
\fancyfoot[C]{\thepage}
\renewcommand{\headrulewidth}{0.4pt}

% ── Mathematics ─────────────────────────────────────────────────────────────
\usepackage{amsmath,amssymb,amsthm}
\usepackage{braket}
\usepackage{mathtools}

% ── Figures & tables ────────────────────────────────────────────────────────
\usepackage{graphicx}
\usepackage{booktabs}
\usepackage{array}
\usepackage{float}
\usepackage{caption}
\captionsetup{font=small,labelfont=bf}
\usepackage{enumitem}

% ── Coloured boxes ──────────────────────────────────────────────────────────
\usepackage{xcolor}
\usepackage[most]{tcolorbox}

\definecolor{conceptbg}{HTML}{EEF2FF}
\definecolor{conceptframe}{HTML}{4338CA}
\definecolor{intuitionbg}{HTML}{ECFDF5}
\definecolor{intuitionframe}{HTML}{059669}
\definecolor{warningbg}{HTML}{FEF2F2}
\definecolor{warningframe}{HTML}{DC2626}
\definecolor{examplebg}{HTML}{FFFBEB}
\definecolor{exampleframe}{HTML}{D97706}
\definecolor{notebookbg}{HTML}{F5F0FF}
\definecolor{notebookframe}{HTML}{7C4DFF}

\newtcolorbox{concept}[1][]{%
  colback=conceptbg, colframe=conceptframe,
  fonttitle=\bfseries, title={Key Concept}, boxrule=0.6pt,
  arc=3pt, left=8pt, right=8pt, top=6pt, bottom=6pt, #1}

\newtcolorbox{intuition}[1][]{%
  colback=intuitionbg, colframe=intuitionframe,
  fonttitle=\bfseries, title={Intuition}, boxrule=0.6pt,
  arc=3pt, left=8pt, right=8pt, top=6pt, bottom=6pt, #1}

\newtcolorbox{warning}[1][]{%
  colback=warningbg, colframe=warningframe,
  fonttitle=\bfseries, title={Subtlety}, boxrule=0.6pt,
  arc=3pt, left=8pt, right=8pt, top=6pt, bottom=6pt, #1}

\newtcolorbox{example}[1][]{%
  colback=examplebg, colframe=exampleframe,
  fonttitle=\bfseries, title={Worked Example}, boxrule=0.6pt,
  arc=3pt, left=8pt, right=8pt, top=6pt, bottom=6pt, #1}

\newtcolorbox{notebook}[1][]{%
  colback=notebookbg, colframe=notebookframe,
  fonttitle=\bfseries, boxrule=0.6pt,
  arc=3pt, left=8pt, right=8pt, top=6pt, bottom=6pt, #1}

% ── Theorem environments ───────────────────────────────────────────────────
\theoremstyle{definition}
\newtheorem{definition}{Definition}[chapter]
\newtheorem{theorem}{Theorem}[chapter]
\newtheorem{proposition}{Proposition}[chapter]

% ── Cross-references & links ──────────────────────────────────────────────
\usepackage[colorlinks=true,linkcolor=conceptframe,citecolor=conceptframe,urlcolor=conceptframe]{hyperref}
\usepackage{cleveref}

% ── Shorthand ──────────────────────────────────────────────────────────────
\newcommand{\expect}[1]{\langle #1 \rangle}
\newcommand{\code}[1]{\texttt{#1}}
\newcommand{\HH}{\mathcal{H}}
\newcommand{\CC}{\mathcal{C}}
\newcommand{\ketT}{\ket{T}}

% ── Title ──────────────────────────────────────────────────────────────────
\title{%
  \vspace{-1cm}
  {\LARGE\bfseries Companion Compendium}\\[0.6em]
  {\Large Encoded Magic-State Preparation\\
  on the $[\![4,2,2]\!]$ Code}\\[1em]
  {\large A Gentle Introduction to the Physics, Engineering,\\
  and Optimisation Behind the Notebook Series}
}
\author{%
  Companion to the \textsc{autoresearch-quantum} project\\[0.3em]
  {\small\url{https://github.com/saymrwulf/autoresearch-quantum}}
}
\date{April 2026}

\begin{document}
\maketitle
\thispagestyle{empty}

\vfill
\begin{center}
\begin{minipage}{0.85\textwidth}
\small\itshape
This compendium is the ``course textbook'' for the twelve Jupyter notebooks
(across four learning plans) in the \textsc{autoresearch-quantum} project.
Start at \texttt{00\_START\_HERE.ipynb} to choose your plan. This document
is designed to be read before, during, or after working through the
notebooks. Every concept exercised in the notebooks is explained here
with the depth and context that a tutorial session cannot provide. No
prior knowledge of quantum error correction is assumed; familiarity with
linear algebra and complex numbers is helpful.
\end{minipage}
\end{center}
\vspace{2cm}

\tableofcontents

% ========================================================================
\chapter{Why This Project Exists}
\label{ch:motivation}
% ========================================================================

\section{The Promise of Quantum Computing}

A quantum computer manipulates quantum bits---\emph{qubits}---that can
exist in superpositions of $\ket{0}$ and $\ket{1}$. Where a classical
bit is either 0 or 1, a qubit is described by
\begin{equation}
  \ket{\psi} = \alpha\ket{0} + \beta\ket{1},
  \qquad |\alpha|^2 + |\beta|^2 = 1,
\end{equation}
where $\alpha$ and $\beta$ are complex numbers called \emph{amplitudes}.
The constraint $|\alpha|^2 + |\beta|^2 = 1$ ensures that the
probabilities of measuring $\ket{0}$ or $\ket{1}$ sum to one.

Multiple qubits can be \emph{entangled}, meaning the state of one
depends on the state of the others in a way that has no classical
analogue. Entanglement is the engine that gives quantum algorithms their
power: Shor's algorithm for factoring, Grover's algorithm for search,
and quantum simulation of molecules all exploit entanglement.

\section{The Noise Problem}

Today's quantum processors are \emph{noisy}. Every gate operation, every
idle moment, and every measurement introduces errors. A typical
two-qubit gate on current IBM hardware has an error rate of roughly 1\%,
which sounds small but compounds rapidly: a circuit of 100 two-qubit
gates has only about a 37\% chance of executing perfectly
($0.99^{100} \approx 0.37$).

This is the central engineering challenge of our era in quantum
computing: the algorithms we want to run require thousands of gates, but
the hardware can barely manage a few hundred before errors dominate.

\section{Quantum Error Correction: The Path Forward}

The solution is \emph{quantum error correction} (QEC). The idea is
conceptually simple: spread the information of one \emph{logical qubit}
across several \emph{physical qubits}, in such a way that errors can be
detected and corrected without disturbing the encoded information.

This is analogous to classical error-correcting codes (like the parity
checks in your hard drive), but with a quantum twist: you cannot copy
a qubit (the \emph{no-cloning theorem}), and measurement generally
destroys the state. Quantum codes must work around both constraints.

\section{The Role of Magic States}

Even with error correction, there is a fundamental limitation.
The \emph{Eastin--Knill theorem} states:

\begin{concept}[title={The Eastin--Knill Theorem}]
No quantum error-correcting code admits a universal set of
\emph{transversal} gates. That is, you cannot implement every gate you
need by simply applying the same operation to each physical qubit
independently.
\end{concept}

Most codes can implement Clifford gates (Hadamard $H$, phase gate $S$,
CNOT) transversally, but Cliffords alone are not enough. The
\emph{Gottesman--Knill theorem} proves that any circuit built entirely
from Clifford gates can be efficiently simulated on a classical
computer---no quantum advantage.

To break out of this trap, you need a \emph{non-Clifford resource}. The
simplest and most common is the $T$~gate (a $\pi/8$ rotation). Rather
than applying $T$ directly on the encoded qubits (which would not be
transversal), the standard approach is:

\begin{enumerate}
  \item Prepare a special auxiliary state called the \textbf{magic state}
        $\ketT$.
  \item Consume $\ketT$ via \emph{gate teleportation} to apply the
        $T$~gate to an encoded qubit.
\end{enumerate}

This compendium---and the entire \textsc{autoresearch-quantum}
project---is about step~1: preparing magic states as reliably and
efficiently as possible.

\section{What This Project Does}

The project builds a complete pipeline for encoded magic-state
preparation on the $[\![4,2,2]\!]$ quantum error-detecting code:

\begin{enumerate}
  \item \textbf{Prepare} the magic state $\ketT$ using different gate
        sequences (seed styles).
  \item \textbf{Encode} it into the $[\![4,2,2]\!]$ code.
  \item \textbf{Verify} the encoding using stabiliser measurements.
  \item \textbf{Measure} the quality using a magic-state witness.
  \item \textbf{Score} the experiment balancing quality, acceptance rate,
        and cost.
  \item \textbf{Optimise} the parameters automatically using a ratchet
        that learns from its own results.
\end{enumerate}

The notebooks let you see, interact with, and modify every step.

% ========================================================================
\chapter{Qubits, Gates, and Circuits}
\label{ch:qubits}
% ========================================================================

\section{The Qubit}

A single qubit lives in a two-dimensional complex vector space
$\HH = \mathbb{C}^2$, with the \emph{computational basis}
$\{\ket{0}, \ket{1}\}$. Any pure state is
\begin{equation}
  \ket{\psi} = \alpha\ket{0} + \beta\ket{1}
  = \begin{pmatrix} \alpha \\ \beta \end{pmatrix},
  \qquad |\alpha|^2 + |\beta|^2 = 1.
\end{equation}

\subsection{The Bloch Sphere}

Every single-qubit state can be visualised as a point on the
\emph{Bloch sphere}. Writing
$\ket{\psi} = \cos(\theta/2)\ket{0} + e^{i\phi}\sin(\theta/2)\ket{1}$,
the state maps to the point $(\sin\theta\cos\phi,\;
\sin\theta\sin\phi,\; \cos\theta)$ on a unit sphere.

\begin{itemize}
  \item $\ket{0}$ is at the north pole $(0,0,1)$.
  \item $\ket{1}$ is at the south pole $(0,0,-1)$.
  \item States on the equator ($\theta = \pi/2$) have equal probability
        of being measured as $\ket{0}$ or $\ket{1}$.
\end{itemize}

\subsection{Global Phase}

If you multiply the entire state by $e^{i\gamma}$, you get a new vector
$e^{i\gamma}\ket{\psi}$ that \emph{cannot be distinguished from
$\ket{\psi}$ by any measurement}. This factor is called the
\emph{global phase}, and it is physically irrelevant.

\begin{warning}
Two states that differ only by a global phase are the same physical
state. Their fidelity is 1.0, and they occupy the same point on the
Bloch sphere. In the notebooks, you will see three seed styles that
produce different-looking amplitude vectors but fidelity~1.0---this is
why.
\end{warning}

\section{Quantum Gates}

Quantum gates are \emph{unitary} transformations: operations that
preserve the norm of the state vector. Every gate $U$ satisfies
$U^\dagger U = I$.

\subsection{Single-Qubit Gates}

The most important single-qubit gates:

\begin{center}
\renewcommand{\arraystretch}{1.3}
\begin{tabular}{lll}
\toprule
\textbf{Gate} & \textbf{Matrix} & \textbf{What it does} \\
\midrule
Pauli $X$ & $\begin{psmallmatrix} 0&1\\1&0 \end{psmallmatrix}$ &
  Bit-flip: $\ket{0}\leftrightarrow\ket{1}$ \\
Pauli $Z$ & $\begin{psmallmatrix} 1&0\\0&-1 \end{psmallmatrix}$ &
  Phase-flip: $\ket{1}\mapsto -\ket{1}$ \\
Pauli $Y$ & $\begin{psmallmatrix} 0&-i\\i&0 \end{psmallmatrix}$ &
  Both bit-flip and phase-flip: $Y = iXZ$ \\
Hadamard $H$ & $\frac{1}{\sqrt{2}}\begin{psmallmatrix} 1&1\\1&-1 \end{psmallmatrix}$ &
  Creates superposition: $\ket{0}\mapsto\ket{+}$ \\
Phase $S$ & $\begin{psmallmatrix} 1&0\\0&i \end{psmallmatrix}$ &
  Quarter-turn around $Z$ \\
$T$ gate & $\begin{psmallmatrix} 1&0\\0&e^{i\pi/4} \end{psmallmatrix}$ &
  Eighth-turn around $Z$: the key non-Clifford gate \\
\bottomrule
\end{tabular}
\end{center}

The gates $\{H, S, \text{CNOT}\}$ generate the \emph{Clifford group}.
Adding $T$ promotes the set to a \emph{universal} gate set: any unitary
can be approximated to arbitrary precision.

\subsection{Two-Qubit Gates}

The most important two-qubit gate is the \textbf{CNOT} (controlled-NOT,
also called CX):
\begin{equation}
  \text{CNOT} = \ket{0}\bra{0}\otimes I + \ket{1}\bra{1}\otimes X
  = \begin{pmatrix} 1&0&0&0\\ 0&1&0&0\\ 0&0&0&1\\ 0&0&1&0 \end{pmatrix}.
\end{equation}
It flips the target qubit if and only if the control qubit is
$\ket{1}$. CNOT is the primary entangling gate and the dominant source
of noise on current hardware.

\begin{intuition}
On IBM hardware, two-qubit gates have error rates 10--100$\times$
higher than single-qubit gates. Minimising the two-qubit gate count is
the single most impactful optimisation for circuit quality.
\end{intuition}

\section{Circuits}

A quantum circuit is a sequence of gates applied to a register of
qubits. Time flows left to right. Qubits are drawn as horizontal lines
(``wires''). Each gate is a box or symbol on its wire(s).

After the gates, you \emph{measure} some or all qubits, collapsing
their superposition into classical bits. A \emph{shot} is one execution
of the full circuit (preparation + gates + measurement).

% ========================================================================
\chapter{The Magic State}
\label{ch:magic}
% ========================================================================

\section{Definition}

The magic state is defined as:
\begin{equation}
  \ketT = \frac{\ket{0} + e^{i\pi/4}\ket{1}}{\sqrt{2}}.
  \label{eq:tstate}
\end{equation}

\begin{concept}[title={T-state Properties}]
\begin{itemize}[nosep]
  \item Amplitudes: $\alpha = 1/\sqrt{2}$,
        $\beta = e^{i\pi/4}/\sqrt{2} = (1+i)/(2)$.
  \item Phase of $\ket{1}$ coefficient: $\pi/4 = 45^\circ$.
  \item Bloch sphere: on the equator ($\expect{Z} = 0$), at $45^\circ$
        between the $+X$ and $+Y$ axes.
  \item $\expect{X} = \expect{Y} = 1/\sqrt{2} \approx 0.7071$.
\end{itemize}
\end{concept}

\begin{warning}[title={$\pi/4$ vs.\ $\pi/8$}]
The gate is called ``$T$'' and sometimes the ``$\pi/8$~gate'' because
of Bloch-sphere conventions (the rotation \emph{angle} is $\pi/4$, but
the \emph{half-angle} in the Bloch parametrisation is $\pi/8$). The
\emph{state} has phase $\pi/4$. The notebooks use $\pi/4$ consistently.
\end{warning}

\section{Why Magic States Matter}

As discussed in \cref{ch:motivation}, the Clifford group alone is
classically simulable (Gottesman--Knill theorem). The $T$~gate breaks
this barrier. But implementing $T$ transversally on most error-correcting
codes is impossible (Eastin--Knill). The workaround:

\begin{enumerate}
  \item Prepare $\ketT$ in an auxiliary register.
  \item Use \emph{gate teleportation}: a circuit of Clifford gates plus
        a measurement that effectively applies $T$ to the target qubit,
        consuming $\ketT$ in the process.
  \item If the $\ketT$ is noisy, apply \emph{magic-state distillation}
        to purify it (at the cost of more copies).
\end{enumerate}

Our project focuses on step~1: preparing the highest-quality $\ketT$
we can, encoded in an error-detecting code, so that downstream
distillation (if needed) starts from the best possible input.

\section{Three Ways to Prepare $\ketT$}
\label{sec:seeds}

On a single qubit, the magic state can be prepared by several
equivalent gate sequences. We call these \emph{seed styles}:

\begin{center}
\renewcommand{\arraystretch}{1.3}
\begin{tabular}{lll}
\toprule
\textbf{Style} & \textbf{Gates} & \textbf{Notes} \\
\midrule
\code{h\_p} & $H$ then $P(\pi/4)$ &
  Most natural: Hadamard creates $\ket{+}$, phase gate adds $\pi/4$ \\
\code{ry\_rz} & $R_Y(\pi/2)$ then $R_Z(\pi/4)$ &
  Native on many hardware platforms \\
\code{u\_magic} & $U(\pi/2, \pi/4, 0)$ &
  Single parameterised gate \\
\bottomrule
\end{tabular}
\end{center}

All three produce the same physical state (fidelity~$= 1.0$). The
amplitude vectors may look different because they differ by a
\emph{global phase}---which, as we discussed, is unphysical. The choice
of seed style matters only when the circuit is \emph{transpiled} for a
specific hardware backend, because different decompositions lead to
different native-gate counts and thus different noise profiles.

% ========================================================================
\chapter{The $[\![4,2,2]\!]$ Error-Detecting Code}
\label{ch:code}
% ========================================================================

\section{Why Encode?}

A bare qubit has no protection against errors. If a cosmic ray or a
stray photon flips a qubit, the computation is silently corrupted. We
need a way to detect (and ideally correct) such errors.

Classical error detection is straightforward: store redundant copies and
compare them. But the \emph{no-cloning theorem} forbids copying an
unknown quantum state:

\begin{theorem}[No-Cloning]
There is no unitary operation $U$ such that
$U\ket{\psi}\ket{0} = \ket{\psi}\ket{\psi}$ for all $\ket{\psi}$.
\end{theorem}

Quantum error correction circumvents this by encoding information not
in copies but in \emph{entanglement patterns}. The information is spread
across multiple physical qubits in a way that individual errors can be
detected without revealing the encoded data.

\section{Code Parameters}

The $[\![4,2,2]\!]$ code is the smallest quantum error-\emph{detecting}
code. Its parameters mean:

\begin{center}
\renewcommand{\arraystretch}{1.3}
\begin{tabular}{cl}
\toprule
\textbf{Parameter} & \textbf{Meaning} \\
\midrule
$n = 4$ & 4 physical qubits \\
$k = 2$ & 2 logical qubits encoded \\
$d = 2$ & Distance 2: detects any single-qubit error \\
\bottomrule
\end{tabular}
\end{center}

\begin{concept}[title={Detection vs.\ Correction}]
Distance $d = 2$ means the code can \emph{detect} any error affecting
a single qubit, but it \emph{cannot correct} it. If an error is
detected, the shot is discarded (postselection). A code needs
distance $d \geq 3$ to correct single-qubit errors.
\end{concept}

For our purposes, detection is sufficient: we discard corrupted shots
and keep only clean ones. This trades quantity for quality---a deliberate
choice that the scoring formula captures.

\section{The Codespace}

The 4~physical qubits span a Hilbert space of dimension $2^4 = 16$.
The $[\![4,2,2]\!]$ code selects a 4-dimensional subspace called the
\emph{codespace} $\CC$. Within $\CC$, the two logical qubits can be in
any state---giving us $2^2 = 4$ degrees of freedom, as expected.

The codespace is defined by two \emph{stabiliser} operators:

\begin{definition}[Stabilisers of the $[\![4,2,2]\!]$ code]
\begin{equation}
  S_X = X \otimes X \otimes X \otimes X = XXXX, \qquad
  S_Z = Z \otimes Z \otimes Z \otimes Z = ZZZZ.
\end{equation}
A state $\ket{\psi}$ is in the codespace if and only if
$S_X\ket{\psi} = +\ket{\psi}$ and $S_Z\ket{\psi} = +\ket{\psi}$.
\end{definition}

In other words, the codespace is the simultaneous $+1$ eigenspace of
both stabilisers. Any state outside this eigenspace has been corrupted
by an error.

\subsection{Properties of the Stabilisers}

Both stabilisers have important algebraic properties:

\begin{enumerate}
  \item \textbf{Squaring to identity:} $S_X^2 = I$ and $S_Z^2 = I$.
        Since $S^2 = I$, the eigenvalues of $S$ can only be $\pm 1$.
  \item \textbf{Commutation:} $[S_X, S_Z] = S_X S_Z - S_Z S_X = 0$.
        The two stabilisers commute, so they can be measured
        simultaneously (they share a common eigenbasis).
  \item \textbf{Hermiticity:} Both are Hermitian ($S^\dagger = S$), so
        they are valid observables.
\end{enumerate}

\section{Logical Operators}
\label{sec:logical-ops}

Within the codespace, we need operators that act on the \emph{logical}
qubits without leaving the codespace. These must commute with both
stabilisers.

For our encoded magic state, the relevant logical operators are:

\begin{center}
\renewcommand{\arraystretch}{1.3}
\begin{tabular}{llll}
\toprule
\textbf{Operator} & \textbf{Pauli string} & \textbf{Qubits acted on} & \textbf{Role} \\
\midrule
Logical $X$ & $IXIX$ & 0, 2 & $X$ on the magic logical qubit \\
Logical $Y$ & $IXZY$ & 0, 1, 2 & $Y$ on the magic logical qubit \\
Spectator $Z$ & $ZIZI$ & 1, 3 & $Z$ on the spectator logical qubit \\
\bottomrule
\end{tabular}
\end{center}

\begin{intuition}
Why does logical $Y$ act on 3 physical qubits? Because the logical
information is \emph{distributed} across all physical qubits by the
encoding. Logical operators must act on this distributed encoding.
There is no single ``logical qubit wire'' to put a $Y$~gate on.
\end{intuition}

\subsection{The Two Logical Qubits}

The $[\![4,2,2]\!]$ code encodes two logical qubits:
\begin{enumerate}
  \item \textbf{Logical qubit 0 (``the magic qubit''):} Prepared in the
        magic state $\ketT$. We measure $\expect{X_L}$ and $\expect{Y_L}$
        to assess its quality.
  \item \textbf{Logical qubit 1 (``the spectator''):} Prepared in
        $\ket{0}_L$. We measure $\expect{Z_{\text{spec}}}$ to confirm it
        has not been disturbed. Ideally $\expect{Z_{\text{spec}}} = +1$.
\end{enumerate}

\section{The Encoder Circuit}
\label{sec:encoder}

The encoder is a unitary circuit that maps a product state on 4~qubits
into an entangled codeword:
\begin{equation}
  \ket{\psi}_L = U_{\text{enc}} \bigl(\ket{T}_0 \otimes \ket{0}_1
  \otimes \ket{0}_2 \otimes \ket{0}_3\bigr).
\end{equation}

The project implements two encoder styles:

\paragraph{cx\_chain (5 CNOT gates, depth 7):} A cascade of CNOT gates
that entangles all four qubits. The Hadamard on qubit~3 creates the
necessary superposition for the second logical qubit (the spectator).
This is the default encoder.

\paragraph{cz\_compiled (5 CZ gates, depth 11):} Uses controlled-$Z$
(CZ) gates instead of CNOT. CZ is the native two-qubit gate on some
hardware platforms (e.g.\ Google's processors). Although the depth is
higher, transpilation may produce fewer native gates on CZ-native
hardware.

Both encoders produce the same logical state (fidelity~$= 1.0$).
The choice between them is an \emph{engineering} decision: which
decomposition produces fewer errors after transpilation for a specific
backend?

\subsection{The Encoded State}

After encoding with the default \code{cx\_chain}, the magic state has
4~non-zero amplitudes out of 16 possible basis states:
\begin{equation}
  \ket{T}_L = \tfrac{1}{2}\ket{0000}
  + \tfrac{e^{i\pi/4}}{2}\ket{0101}
  + \tfrac{e^{i\pi/4}}{2}\ket{1010}
  + \tfrac{1}{2}\ket{1111}.
\end{equation}

The pattern $\{0000, 0101, 1010, 1111\}$ is the codespace of the
$[\![4,2,2]\!]$ code. Any measurement outcome outside these four
bitstrings indicates an error.

\section{Error Detection}
\label{sec:errors}

The power of the code lies in its ability to detect errors. Consider
what happens when a single-qubit Pauli error acts on the encoded state:

\begin{center}
\renewcommand{\arraystretch}{1.3}
\begin{tabular}{lccc}
\toprule
\textbf{Error type} & $\expect{XXXX}$ & $\expect{ZZZZ}$ &
\textbf{Detected by} \\
\midrule
No error & $+1$ & $+1$ & --- \\
$X$ on any qubit & $+1$ & $-1$ & $ZZZZ$ \\
$Z$ on any qubit & $-1$ & $+1$ & $XXXX$ \\
$Y$ on any qubit & $-1$ & $-1$ & Both \\
\bottomrule
\end{tabular}
\end{center}

\begin{concept}[title={Why ZZZZ Detects X Errors}]
$ZZZZ$ detects $X$~errors because $X$ and $Z$ \emph{anti-commute}:
$XZ = -ZX$. When an $X$~error occurs on one qubit, it flips the sign
of the $ZZZZ$ eigenvalue from $+1$ to $-1$.

Similarly, $XXXX$ detects $Z$~errors. A $Y$~error ($Y = iXZ$) triggers
both stabilisers.
\end{concept}

Every single-qubit error flips at least one stabiliser---this is the
\emph{distance-2 guarantee}. A weight-2 error (two qubits affected
simultaneously) could go undetected; that is the inherent limitation of
distance~2.

% ========================================================================
\chapter{Measurement, Verification, and Postselection}
\label{ch:measurement}
% ========================================================================

\section{The Measurement Problem}

To check whether an error has occurred, we need to measure the
stabilisers $\expect{XXXX}$ and $\expect{ZZZZ}$. But there is a
problem: \emph{directly measuring the data qubits collapses the
superposition and destroys the encoded state}.

\section{Ancilla-Based Syndrome Extraction}

The solution is to use \emph{ancilla qubits}---extra qubits that are
entangled with the data qubits in a controlled way, then measured. The
measurement outcome of the ancilla (called the \emph{syndrome}) tells us
whether an error occurred, without revealing the actual data.

\begin{intuition}
Think of the ancilla as a thermometer: it reads the ``temperature'' of
the data qubits (error or no error) without disturbing the patient
(the encoded state). The syndrome bits are the thermometer's reading.
\end{intuition}

In the $[\![4,2,2]\!]$ code, we extract two syndrome bits:
\begin{itemize}
  \item Syndrome bit for $ZZZZ$: $+1$ (no $X$~error) or $-1$ ($X$~error
        detected).
  \item Syndrome bit for $XXXX$: $+1$ (no $Z$~error) or $-1$ ($Z$~error
        detected).
\end{itemize}

We encode these as bits: 0 means $+1$ (no error), 1 means $-1$ (error
detected). The syndrome string ``00'' means both stabilisers are
satisfied---the state is in the codespace.

\section{Postselection}
\label{sec:postselection}

\emph{Postselection} is the process of keeping only the shots where the
syndrome indicates no error:

\begin{definition}[Postselection]
Given $N$ total shots, let $N_{\text{accept}}$ be the number with
syndrome ``00''. The \emph{acceptance rate} is
\begin{equation}
  r_{\text{accept}} = \frac{N_{\text{accept}}}{N}.
\end{equation}
All other shots are discarded.
\end{definition}

\begin{concept}[title={The Cost of Postselection}]
Postselection improves quality by filtering out corrupted shots. But it
\emph{reduces the number of usable data points}. If the acceptance rate
is 60\%, you need $\sim\!1.7\times$ as many total shots to get the same
statistical power. The scoring formula accounts for this trade-off.
\end{concept}

On an ideal (noiseless) simulator, the acceptance rate is 100\%---every
shot passes the syndrome check. On a noisy backend, some shots will
trigger the syndrome flag and be discarded. The acceptance rate is
therefore a direct measure of how much noise affects the experiment.

\section{Witness Circuits}
\label{sec:witness-circuits}

After postselection, we need to measure the logical operators to assess
the quality of the encoded magic state. This requires three separate
circuits:

\begin{enumerate}
  \item \textbf{Logical $X$ circuit:} Measures $\expect{X_L}$ by
        measuring qubits 0 and 2 in the $X$ basis.
  \item \textbf{Logical $Y$ circuit:} Measures $\expect{Y_L}$ by
        measuring qubits 0, 1, 2 in the appropriate bases.
  \item \textbf{Spectator $Z$ circuit:} Measures $\expect{Z_{\text{spec}}}$
        by measuring qubits 1 and 3 in the $Z$ basis.
\end{enumerate}

\begin{warning}[title={Why Three Separate Circuits?}]
The three logical operators \emph{do not commute} with each other.
Measuring one would disturb the quantum state in a way that invalidates
the measurement of the others. Therefore, each operator requires its own
copy of the experiment.
\end{warning}

% ========================================================================
\chapter{The Magic-State Witness}
\label{ch:witness}
% ========================================================================

\section{What Is a Witness?}

A \emph{witness} is a number computed from measurement results that
quantifies how close the prepared state is to the ideal target.
Unlike fidelity (which requires full state tomography), the witness
uses only a few expectation values and can be estimated efficiently.

\section{The Formula}
\label{sec:witness-formula}

The magic-state witness for the $[\![4,2,2]\!]$ encoded $\ketT$ is:

\begin{equation}
  \boxed{
  W = \underbrace{\frac{1 + \frac{\expect{X_L} + \expect{Y_L}}{\sqrt{2}}}{2}}_{\text{magic factor}}
  \;\times\;
  \underbrace{\frac{1 + \expect{Z_{\text{spec}}}}{2}}_{\text{spectator factor}}
  }
  \label{eq:witness}
\end{equation}

\subsection{The Magic Factor}

The magic factor measures how well the encoded state matches the
$T$-state character. For the ideal $\ketT$:
\begin{itemize}
  \item $\expect{X_L} = 1/\sqrt{2} \approx 0.7071$
  \item $\expect{Y_L} = 1/\sqrt{2} \approx 0.7071$
  \item Magic factor $= (1 + (0.7071 + 0.7071)/\sqrt{2})/2 = (1+1)/2 = 1.0$
\end{itemize}

\subsection{The Spectator Factor}

The spectator factor checks that the second logical qubit (which should
be in $\ket{0}_L$) has not been disturbed:
\begin{itemize}
  \item Ideal: $\expect{Z_{\text{spec}}} = +1$
  \item Spectator factor $= (1+1)/2 = 1.0$
\end{itemize}

\subsection{Ideal Witness Value}

For a perfect preparation:
\begin{equation}
  W_{\text{ideal}} = 1.0 \times 1.0 = 1.0.
\end{equation}

Any noise or error reduces $W$ below 1. The witness is deliberately
\emph{sensitive}: even moderate noise produces a noticeable drop, making
it a useful diagnostic.

\begin{example}[title={Witness Under Noise}]
Suppose noise reduces the logical expectations to $\expect{X_L} = 0.55$,
$\expect{Y_L} = 0.50$, and $\expect{Z_{\text{spec}}} = 0.02$:
\begin{align}
  \text{Magic factor} &= \frac{1 + (0.55 + 0.50)/\sqrt{2}}{2}
                       = \frac{1 + 0.742}{2} = 0.871, \\
  \text{Spectator factor} &= \frac{1 + 0.02}{2} = 0.510, \\
  W &= 0.871 \times 0.510 = 0.444.
\end{align}
The witness dropped from 1.0 to 0.44---a clear signal that quality has
degraded significantly.
\end{example}

\section{Witness vs.\ Fidelity}

\emph{Fidelity} measures the overlap between the actual state and the
ideal state: $F = |\!\braket{T_L|\rho|T_L}\!|$. It requires full
knowledge of the density matrix (state tomography), which is expensive.

The \emph{witness} is a proxy that can be estimated from just three
expectation values. It is not identical to fidelity, but it tracks
quality faithfully and is operationally efficient.

% ========================================================================
\chapter{Noise and the Hardware Reality}
\label{ch:noise}
% ========================================================================

\section{Sources of Noise}

On real quantum hardware, errors arise from:

\begin{enumerate}
  \item \textbf{Gate errors:} Imperfect implementation of unitary
        operations. Two-qubit gates are the worst offenders, with error
        rates of $10^{-3}$ to $10^{-2}$.
  \item \textbf{Readout errors:} Misidentification of $\ket{0}$ vs.\
        $\ket{1}$ during measurement.
  \item \textbf{Decoherence:} The qubit loses its quantum properties
        over time ($T_1$ relaxation and $T_2$ dephasing).
  \item \textbf{Cross-talk:} Operations on one qubit inadvertently
        affect neighbouring qubits.
\end{enumerate}

\section{Noise Models and Simulators}

Since access to real quantum hardware is limited and expensive, we use
\emph{noise-model simulators}. The project uses Qiskit Aer's
\code{AerSimulator} with noise models extracted from real IBM backends:

\begin{itemize}
  \item \code{fake\_brisbane}: A 127-qubit noise model mimicking the IBM
        Brisbane processor, with realistic gate error rates, readout
        errors, and $T_1$/$T_2$ times.
\end{itemize}

\begin{intuition}
A noise-model simulator is like a flight simulator: it reproduces the
conditions of the real thing (including turbulence) without the risk or
cost of actual flight time. Results are statistically realistic, even
though no quantum hardware is involved.
\end{intuition}

\section{Transpilation: From Logical to Physical}
\label{sec:transpilation}

The circuit you write in Qiskit uses abstract gates like $H$, CNOT, and
$T$. But real hardware only supports a limited set of \emph{native gates}
(e.g.\ IBM's basis: $\{$ECR, $R_Z$, $S_X$, $X$$\}$). The
\emph{transpiler} converts your circuit into native gates:

\begin{enumerate}
  \item \textbf{Gate decomposition:} $H \to S_X \cdot R_Z(\pi/2)$, etc.
  \item \textbf{Qubit routing:} Map logical qubits to physical qubits on
        the hardware's connectivity graph. Insert SWAP gates where
        needed.
  \item \textbf{Optimisation:} Cancel redundant gates, merge rotations,
        simplify sequences.
\end{enumerate}

Qiskit provides optimisation levels 0--3:

\begin{center}
\renewcommand{\arraystretch}{1.3}
\begin{tabular}{cl}
\toprule
\textbf{Level} & \textbf{Description} \\
\midrule
0 & No optimisation (just decomposition and routing) \\
1 & Light optimisation (default) \\
2 & Medium: gate cancellation, commutation analysis \\
3 & Heavy: resynthesis of 2-qubit blocks \\
\bottomrule
\end{tabular}
\end{center}

\begin{warning}[title={Higher Is Not Always Better}]
Aggressive optimisation reduces gate count but may reroute qubits onto
noisier connections. The \emph{net} effect depends on the specific
circuit and the specific hardware topology. The ratchet explores multiple
optimisation levels to find the empirically best choice.
\end{warning}

\section{Cost Model}

The \emph{cost} of a circuit quantifies its resource consumption. The
project uses a weighted sum:
\begin{equation}
  \text{cost} = w_{\text{2q}} \cdot n_{\text{2q}}
              + w_{\text{depth}} \cdot d + c_0,
\end{equation}
where $n_{\text{2q}}$ is the two-qubit gate count, $d$ is the circuit
depth, and $c_0$ is a baseline cost. The weights come from the rung
configuration and can be tuned.

Two-qubit gates dominate: they are the noisiest operations, so
$w_{\text{2q}}$ is typically the largest weight.

% ========================================================================
\chapter{Scoring: Putting It All Together}
\label{ch:scoring}
% ========================================================================

An experiment produces several metrics: witness value, acceptance rate,
circuit cost. We need a single number to compare experiments. The
\emph{score} does this.

\section{The Weighted Acceptance-Cost Score}

For rungs 1--3, the score is:
\begin{equation}
  \boxed{
    \text{score} = \frac{\text{quality} \times \text{acceptance\_rate}}{\text{cost}}
  }
  \label{eq:wac}
\end{equation}
where \emph{quality} is the magic-state witness $W$.

\begin{intuition}
The score answers: ``How much magic-state quality do I get per unit of
resource spent, accounting for the shots I had to throw away?'' It
naturally balances three competing concerns:
\begin{itemize}[nosep]
  \item \textbf{Quality:} higher witness is better.
  \item \textbf{Acceptance:} fewer discarded shots is better.
  \item \textbf{Cost:} simpler circuits are better (cheaper to run).
\end{itemize}
\end{intuition}

\begin{example}[title={Score Tension}]
A circuit that is $2\times$ better in quality but $3\times$ more
expensive scores \emph{worse}:
\[
  \frac{2q \cdot a}{3c} = \frac{2}{3} \cdot \frac{qa}{c}
  < \frac{qa}{c}.
\]
The score penalises complexity unless it delivers proportionally more
quality.
\end{example}

\section{Factory Throughput Score}

For rungs 4--5, the score shifts to a \emph{factory throughput} model:
how many usable $T$-states per unit time can the circuit produce?

\begin{equation}
  \text{factory\_score} = \text{quality} \times \text{acceptance\_rate}
  \times \frac{1}{\text{cost}^{3/2}}.
\end{equation}

The steeper cost penalty ($\text{cost}^{3/2}$ instead of $\text{cost}$)
reflects the operational reality that in a $T$-state factory, cost
compounds: each round of distillation consumes multiple copies.

\section{Failure Modes}
\label{sec:failures}

Three things can go wrong, in increasing order of severity:

\begin{enumerate}
  \item \textbf{High cost:} The circuit is expensive but still works.
        Fix: optimise transpilation settings.
  \item \textbf{Poor acceptance:} Many shots are rejected. This wastes
        compute but the accepted shots may still be good. Fix: reduce
        noise exposure (fewer gates, better layout).
  \item \textbf{Low magic witness:} The $T$-state character itself is
        lost. Even the accepted shots produce poor quality. This is the
        most severe failure---the experiment has failed its fundamental
        purpose.
\end{enumerate}

% ========================================================================
\chapter{The Ratchet: Learning by Doing}
\label{ch:ratchet}
% ========================================================================

The ratchet is an automated optimisation system inspired by Andrej
Karpathy's ``autoresearch'' philosophy: let the system run experiments,
learn from the results, and improve its own configuration.

\section{The Incumbent-Challenger Model}

\begin{concept}[title={The Ratchet Guarantee}]
The ratchet maintains an \emph{incumbent}---the best configuration found
so far. In each step, it generates \emph{challengers} (alternative
configurations) and evaluates them. A challenger replaces the incumbent
\emph{only if it scores strictly higher}. The incumbent never gets
worse.
\end{concept}

This monotonicity guarantee is the defining property of a ratchet (named
after the mechanical device that turns in only one direction). It means
the search is safe: you can stop at any time and your best result is
preserved.

\subsection{The Bootstrap Incumbent}

The first incumbent is not random---it is a \emph{bootstrap
incumbent}: a hand-picked, domain-expert guess at reasonable default
parameters. This warm start means the ratchet begins from a sensible
baseline rather than wasting time on obviously bad configurations.

\section{Challenger Generation Strategies}
\label{sec:strategies}

The ratchet generates challengers using three strategies, each with a
budget allocation:

\subsection{NeighborWalk (40\% of budget)}

Changes \emph{exactly one parameter} at a time, trying every alternative
value for that parameter while keeping all others fixed.

\begin{itemize}
  \item \textbf{Strengths:} Systematic, guaranteed to find all
        single-parameter improvements.
  \item \textbf{Weakness:} Blind to \emph{parameter interactions}---it
        cannot discover that changing two parameters simultaneously
        produces a synergy that neither change alone would find.
\end{itemize}

\subsection{RandomCombo (30\% of budget)}

Mutates \emph{multiple parameters simultaneously} at random.

\begin{itemize}
  \item \textbf{Strengths:} Can discover multi-parameter interactions
        and escape local optima.
  \item \textbf{Weakness:} Less systematic---relies on luck to find
        good combinations.
\end{itemize}

\subsection{LessonGuided (30\% of budget)}

Uses rules extracted from previous experiments to focus the search.

\begin{itemize}
  \item \textbf{Strengths:} Exploits accumulated knowledge, avoids
        repeating known mistakes.
  \item \textbf{Weakness:} Only available after lessons have been
        extracted (not on the first rung).
\end{itemize}

\begin{intuition}
Think of the three strategies as three employees searching a warehouse:
NeighborWalk checks one shelf at a time (thorough but slow).
RandomCombo wanders around trying random combinations (creative but
unpredictable). LessonGuided reads the notes from previous searches
first (efficient but needs prior experience).
\end{intuition}

\section{Ratchet Steps and Rungs}

A \textbf{step} is one round of challenger generation and evaluation.
In each step:
\begin{enumerate}
  \item Generate a batch of challengers using the three strategies.
  \item Evaluate each challenger (run the experiment, compute the score).
  \item If any challenger beats the incumbent, replace the incumbent.
  \item Log the result (winner, margin, all scores).
\end{enumerate}

A \textbf{rung} is a sequence of steps, terminated when the
\emph{patience} is exhausted:

\begin{definition}[Patience]
If $p$ consecutive steps fail to improve the incumbent, the rung stops.
This prevents wasting compute once the nearby parameter space has been
exhausted.
\end{definition}

\section{Lesson Extraction}
\label{sec:lessons}

After each rung, the system analyses all experimental results and
extracts \emph{lessons}---rules about which parameter values help or
hurt:

\begin{enumerate}
  \item \textbf{Fix rules:} ``Always use this value''---a parameter
        value that consistently appears in top-scoring experiments.
  \item \textbf{Avoid rules:} ``Never use this value''---a parameter
        value that consistently appears in bottom-scoring experiments.
\end{enumerate}

Lessons are stored in two formats:
\begin{itemize}
  \item A human-readable \emph{narrative} (natural language summary).
  \item Machine-readable \emph{SearchRules} (JSON) that the
        LessonGuided strategy can consume directly.
\end{itemize}

\section{Search Space Narrowing}

Lessons also \emph{narrow} the search space for subsequent rungs.
If a value is consistently bad, it is removed from the allowed options.
The dimension remains (the parameter still exists), but with fewer
values to explore. A minimum number of values per dimension is preserved
to prevent overfitting to noise.

\section{Cross-Rung Propagation}

The winning configuration from rung $N$ is \emph{propagated} as the
bootstrap incumbent for rung $N+1$. This avoids cold-starting each rung
and allows the system to build on previous gains.

Combined with lesson extraction and space narrowing, this creates a
\emph{progressive refinement} loop:
\[
  \text{Run} \to \text{Learn} \to \text{Narrow} \to \text{Propagate}
  \to \text{Run} \to \cdots
\]

\section{Transfer Evaluation}

A configuration optimised for one backend might be overfitted to that
backend's specific noise profile. \emph{Transfer evaluation} tests the
winning configuration on a different backend:

\begin{itemize}
  \item If the score is similar, the configuration is \emph{robust}.
  \item If the score drops sharply, the configuration is
        \emph{overfitted} to the source backend.
\end{itemize}

Transfer evaluation is a form of generalisation testing, analogous to
validating a machine-learning model on a held-out dataset.

% ========================================================================
\chapter{The Five Rungs}
\label{ch:rungs}
% ========================================================================

The system organises its search into five \emph{rungs}, each adding
complexity:

\begin{center}
\renewcommand{\arraystretch}{1.4}
\begin{tabular}{clll}
\toprule
\textbf{Rung} & \textbf{Focus} & \textbf{Backend} & \textbf{Scorer} \\
\midrule
1 & Core parameters (seed, encoder, verification) &
    \code{fake\_brisbane} & WAC \\
2 & Transpilation (opt level, layout, routing) &
    \code{fake\_brisbane} & WAC \\
3 & Fine-tuning (approximation degree, initial layout) &
    \code{fake\_brisbane} & WAC \\
4 & Hardware-aware optimisation &
    \code{fake\_brisbane} & Factory \\
5 & Transfer validation &
    Different backend & Factory \\
\bottomrule
\end{tabular}
\end{center}

Each rung inherits the best configuration from the previous one,
narrows the search space based on lessons, and explores the next layer
of parameters. The first three rungs use the Weighted Acceptance-Cost
(WAC) score; rungs 4--5 switch to factory throughput to reflect
production-oriented priorities.

% ========================================================================
\chapter{Putting It All Together: The Pipeline}
\label{ch:pipeline}
% ========================================================================

Here is the complete flow from start to finish:

\begin{enumerate}
  \item \textbf{Configuration:} Load a rung YAML file specifying the
        parameter space, scorer, backend, and budget.
  \item \textbf{Preparation:} Build the quantum circuit from the
        experiment spec (seed style $\to$ encoder $\to$ verification
        circuits $\to$ witness circuits).
  \item \textbf{Transpilation:} Compile the circuit for the target
        backend at the specified optimisation level.
  \item \textbf{Execution:} Run the circuit for $N$ shots on the
        noise-model simulator.
  \item \textbf{Analysis:}
    \begin{enumerate}[nosep]
      \item Parse syndrome bits to compute acceptance rate.
      \item Parse data bits on accepted shots to compute
            $\expect{X_L}$, $\expect{Y_L}$, $\expect{Z_{\text{spec}}}$.
      \item Compute the witness $W$.
      \item Compute the cost from circuit metrics.
      \item Compute the score.
    \end{enumerate}
  \item \textbf{Ratchet:} Compare challenger scores to the incumbent.
        Promote the winner. Extract lessons. Narrow the search space.
        Propagate to the next rung.
  \item \textbf{Transfer:} Test the final configuration on a different
        backend.
\end{enumerate}

\begin{notebook}[title={Where to See This in the Notebooks}]
\begin{itemize}[nosep]
  \item \textbf{Plan A, Notebook 01:} Steps 1--5 in detail (one cell per
        stage).
  \item \textbf{Plan A, Notebook 02:} Step 5 in depth (scoring,
        parameter sweeps).
  \item \textbf{Plan A, Notebook 03:} Steps 6--7 (the ratchet in
        action).
  \item \textbf{Plan B (Spiral):} All steps in three passes of
        increasing depth.
  \item \textbf{Plan C, Track A:} Steps 1--3 (physics focus).
  \item \textbf{Plan C, Track B:} Steps 3--5 (engineering focus).
  \item \textbf{Plan C, Track C:} Steps 6--7 (optimisation focus).
  \item \textbf{Plan C, Dashboard:} Interactive exploration of step 2
        parameters.
  \item \textbf{Plan D, Experiment 1:} Steps 1--3 (encoding and error
        detection, ideal simulator).
  \item \textbf{Plan D, Experiment 2:} Steps 3--5 (noise, scoring,
        parameter sweep).
  \item \textbf{Plan D, Experiment 3:} Steps 6--7 (ratchet, lessons,
        transfer evaluation).
\end{itemize}
\end{notebook}

% ========================================================================
\chapter{Glossary}
\label{ch:glossary}
% ========================================================================

\begin{description}[style=nextline, leftmargin=3cm, labelwidth=2.8cm]

\item[Acceptance rate]
Fraction of shots that pass the syndrome check (postselection).

\item[Ancilla]
An auxiliary qubit used for syndrome extraction without disturbing the
data qubits.

\item[Bloch sphere]
A geometric representation of a single-qubit state as a point on a unit
sphere.

\item[Bootstrap incumbent]
The hand-picked initial configuration that the ratchet starts from.

\item[Challenger]
A candidate configuration that competes against the incumbent in a
ratchet step.

\item[Clifford group]
The group of gates generated by $\{H, S, \text{CNOT}\}$. Classically
simulable by the Gottesman--Knill theorem.

\item[Codespace]
The subspace of the physical Hilbert space where valid codewords live.
Defined by the simultaneous $+1$ eigenspace of all stabilisers.

\item[Cost]
A scalar measuring the resource consumption of a circuit (dominated by
two-qubit gate count).

\item[Distance]
The minimum weight of an undetectable error. For $[\![4,2,2]\!]$,
$d = 2$: all weight-1 errors are detectable.

\item[Eastin--Knill theorem]
No quantum code admits a universal set of transversal gates.

\item[Factory throughput]
A scoring function that penalises cost more heavily, modelling a
$T$-state production pipeline.

\item[Fidelity]
$F = |\!\braket{\psi|\rho|\psi}\!|$: the overlap between the actual
state and the ideal target.

\item[Global phase]
A factor $e^{i\gamma}$ multiplying the entire state vector. Unphysical
and unmeasurable.

\item[Gottesman--Knill theorem]
Clifford-only circuits can be efficiently simulated on a classical
computer.

\item[Incumbent]
The best configuration found so far. Replaced only when a challenger
scores strictly higher.

\item[Lesson]
A rule extracted from experimental results (``fix'' or ``avoid'' a
parameter value).

\item[Logical operator]
An operator that acts on the encoded (logical) information within the
codespace.

\item[Magic state]
The state $\ketT = (\ket{0} + e^{i\pi/4}\ket{1})/\sqrt{2}$, consumed
to implement the $T$~gate via gate teleportation.

\item[Narrowing]
Removing poorly-performing parameter values from the search space
between rungs.

\item[No-cloning theorem]
No unitary operation can copy an unknown quantum state.

\item[Patience]
The number of consecutive no-improvement steps before a rung terminates.

\item[Postselection]
Discarding shots where the syndrome indicates an error.

\item[Ratchet]
A monotonic optimiser: the incumbent never gets worse.

\item[Rung]
A stage of the optimisation pipeline with specific parameters to explore
and a fixed budget.

\item[Seed style]
The gate sequence used to prepare the magic state on a single qubit
before encoding.

\item[Shot]
One execution of the full circuit (preparation + measurement).

\item[Stabiliser]
A Pauli operator whose $+1$ eigenspace defines the codespace. For
$[\![4,2,2]\!]$: $XXXX$ and $ZZZZ$.

\item[Syndrome]
The measurement outcome of the ancilla qubits, indicating whether an
error has been detected.

\item[Transpilation]
Converting a logical circuit into native gates for a specific hardware
backend.

\item[Transfer evaluation]
Testing a configuration on a different backend to check for overfitting.

\item[Witness]
A scalar quantity computed from expectation values that estimates the
quality of the prepared magic state.

\item[$[\![n,k,d]\!]$]
Standard notation for a quantum code: $n$~physical qubits, $k$~logical
qubits, distance~$d$.

\end{description}

% ========================================================================
\appendix
\chapter{Mathematical Background}
\label{app:math}
% ========================================================================

This appendix collects the mathematical prerequisites. Skip it if you
are comfortable with linear algebra over $\mathbb{C}$.

\section{Complex Numbers and Amplitudes}

A complex number $z = a + bi$ has a real part $a$, imaginary part $b$,
magnitude $|z| = \sqrt{a^2 + b^2}$, and phase
$\arg(z) = \arctan(b/a)$. In polar form: $z = |z| e^{i\theta}$.

Quantum amplitudes are complex numbers. The probability of measuring a
state is the squared magnitude of its amplitude: $p = |\alpha|^2$.

\section{Tensor Products}

The state space of $n$ qubits is the tensor product
$(\mathbb{C}^2)^{\otimes n}$, with dimension $2^n$. For two qubits:
\begin{equation}
  \ket{\psi} \otimes \ket{\phi}
  = \begin{pmatrix} \alpha_\psi \\ \beta_\psi \end{pmatrix}
    \otimes
    \begin{pmatrix} \alpha_\phi \\ \beta_\phi \end{pmatrix}
  = \begin{pmatrix}
      \alpha_\psi \alpha_\phi \\
      \alpha_\psi \beta_\phi \\
      \beta_\psi \alpha_\phi \\
      \beta_\psi \beta_\phi
    \end{pmatrix}.
\end{equation}

\section{Pauli Matrices}

The four Pauli matrices form a basis for $2 \times 2$ Hermitian
matrices:
\begin{equation}
  I = \begin{pmatrix} 1&0\\0&1 \end{pmatrix}, \quad
  X = \begin{pmatrix} 0&1\\1&0 \end{pmatrix}, \quad
  Y = \begin{pmatrix} 0&-i\\i&0 \end{pmatrix}, \quad
  Z = \begin{pmatrix} 1&0\\0&-1 \end{pmatrix}.
\end{equation}

Key properties:
\begin{itemize}
  \item $X^2 = Y^2 = Z^2 = I$.
  \item $XY = iZ$, $YZ = iX$, $ZX = iY$ (cyclic).
  \item $XZ = -ZX$ (anti-commutation). This is why $ZZZZ$ detects
        $X$~errors.
\end{itemize}

A Pauli string like $XXXX$ is the tensor product $X \otimes X \otimes X
\otimes X$, a $16 \times 16$ matrix acting on 4~qubits.

\section{Eigenvalues and Expectation Values}

The \emph{expectation value} of an observable $O$ in state
$\ket{\psi}$ is:
\begin{equation}
  \expect{O} = \bra{\psi} O \ket{\psi}.
\end{equation}

For a Pauli operator $P$ with $P^2 = I$, the eigenvalues are exactly
$\pm 1$. A measurement of $P$ always returns $+1$ or $-1$. The
expectation value is the average over many measurements.

\chapter{Notebook--Compendium Cross-Reference}
\label{app:xref}

\begin{center}
\renewcommand{\arraystretch}{1.3}
\begin{tabular}{p{5.2cm}p{3.6cm}l}
\toprule
\textbf{Notebook Topic} & \textbf{Notebooks} & \textbf{Compendium} \\
\midrule
T-state definition \& Bloch sphere &
  A/01~\S1--2, B~\S2.1, C/A~\S1--3, D/1~\S1 &
  \cref{ch:magic} \\
Why encode (no-cloning, distance) &
  A/01~\S3, C/A~\S1, D/1~\S2 &
  \cref{ch:code}~\S1--2 \\
Stabilisers \& codespace &
  A/01~\S6, B~\S2.3, C/A~\S4, D/1~\S3 &
  \cref{ch:code}~\S3 \\
Logical operators &
  A/01~\S6, C/A~\S5 &
  \cref{sec:logical-ops} \\
Encoder circuits &
  A/01~\S4--5, C/A~\S6 &
  \cref{sec:encoder} \\
Error detection &
  A/01~\S7, C/A~\S8, D/1~\S4 &
  \cref{sec:errors} \\
Ancilla \& syndrome extraction &
  A/01~\S9, C/A~\S7 &
  \cref{ch:measurement}~\S2 \\
Postselection &
  A/01~\S11, A/02~\S3, B~\S2.5, D/1~\S6 &
  \cref{sec:postselection} \\
Noise models \& transpilation &
  A/02~\S2, C/B~\S1--3, D/2~\S1 &
  \cref{ch:noise} \\
Magic witness formula &
  A/02~\S5, B~\S2.7, C/A~\S9, D/1~\S5 &
  \cref{ch:witness} \\
Scoring formula &
  A/02~\S7, B~\S2.9, C/B~\S8, D/2~\S2 &
  \cref{ch:scoring} \\
Factory throughput &
  A/02~\S10, C/B~\S9 &
  \cref{ch:scoring}~\S2 \\
Failure modes &
  A/02~\S9, C/B~\S7 &
  \cref{sec:failures} \\
Ratchet mechanism &
  A/03~\S1--4, B~\S2.10--12, C/C~\S1--7, D/3~\S1--2 &
  \cref{ch:ratchet}~\S1--3 \\
Search strategies &
  A/03~\S7, B~\S3.5, C/C~\S3--4, D/3~\S2 &
  \cref{sec:strategies} \\
Lesson extraction \& rules &
  A/03~\S8, B~\S3.6, C/C~\S8--9, D/3~\S4 &
  \cref{sec:lessons} \\
Narrowing \& propagation &
  B~\S3.7, C/C~\S10--11 &
  \cref{ch:ratchet}~\S5--6 \\
Transfer evaluation &
  A/03~\S10, B~\S3.8, C/C~\S12, D/3~\S5 &
  \cref{ch:ratchet}~\S7 \\
Parameter sweep \& optimisation &
  A/02~\S8, D/2~\S3, D/3~\S3 &
  \cref{ch:scoring} \\
\bottomrule
\end{tabular}
\end{center}

\vfill
\begin{center}
\rule{0.3\textwidth}{0.4pt}\\[0.5em]
{\small End of compendium.}
\end{center}

\end{document}
